\section{Animations}

% \begin{animation}[options] with one \animpanel{label}{path} per panel; \N in
% the path is the frame number (figures/anim/case1_0.png, _2, ... _22). The
% panels play in sync. Options:
%   fps=3  frames=first:last:step  width=  columns=N (rows of N panels)
%   controls=all | none | hover | play,step,stop,speed   autoplay  loop
%   subcaptions ((a), (b) before the labels)
%   dt=0.05 (time = frame x dt below the panels)  timelabel={$t = $}
%   timedigits=2  timeunit=s  timebar (progress bar, clickable in HTML)
% In the PDF the animation plays in Acrobat/Okular; the HTML deck plays it
% everywhere. Needs \usepackage{animate}.
\begin{frame}{Animations}{panels in one row}
  \begin{figure}
    \begin{animation}[fps=4, frames=0:22:2, dt=0.1, timebar, controls=hover]
      \animpanel{Case 1}{figures/anim/case1_\N}
      \animpanel{Case 2}{figures/anim/case2_\N}
      \animpanel{Case 3}{figures/anim/case3_\N}
    \end{animation}
    \caption{Three cases in sync; buttons appear on hover in the HTML deck.}
  \end{figure}
\end{frame}

\begin{frame}{Animations}{panels in a grid: \texttt{columns=2}}
  \begin{columns}[c]
    \begin{column}{0.4\textwidth}
      \begin{itemize}
        \item \texttt{columns=2}: rows of two panels
        \item Each label sits on its panel
        \item \texttt{width=} makes the panels smaller
        \item \texttt{subcaptions}: (a), (b), \ldots
      \end{itemize}
    \end{column}
    \begin{column}{0.56\textwidth}
      \begin{figure}
        \begin{animation}[fps=4, frames=0:22:2, columns=2, width=0.35\linewidth,
                          subcaptions, dt=0.1, timelabel={$\tau = $}]
          \animpanel{Case 1}{figures/anim/case1_\N}
          \animpanel{Case 2}{figures/anim/case2_\N}
          \animpanel{Case 3}{figures/anim/case3_\N}
          \animpanel{Case 4}{figures/anim/case4_\N}
        \end{animation}
        \caption{Four cases, two per row.}
      \end{figure}
    \end{column}
  \end{columns}
\end{frame}
